
\subsection{Measurement protocol}

\paragraph{End-to-end architecture.}
All three end-to-end models are untrained, 1,000-class Tiny variants with a
$4\times4$, stride-4 patch stem; output stride 32; four stages with block depths
$(3,3,9,3)$ and channel widths $(96,192,384,768)$; and a global-average-pooling,
LayerNorm, linear classification head. Blocks use GELU, layer-scale
initialization $10^{-6}$, no convolutional MLP or global response normalization,
and zero dropout and stochastic-depth rates. ConvNeXt-T uses a $7\times7$
depthwise convolution in each block. Both WTConvNeXt-T variants replace it with
$5\times5$ WTConv at stage-wise decomposition levels $(5,4,3,2)$ and differ
only in whether WTConv uses the reference or fused implementation. Measurements
use FP32 inputs of shape $(64,3,224,224)$.

\paragraph{Timing.}
For each layer configuration and end-to-end model, we run $20$ untimed warm-up
iterations followed by $50$ measured iterations. Warm-up excludes first-use
compilation and cuDNN setup. We synchronize CUDA before and after the measured
loop and use CUDA events to record each iteration's device time. Layerwise
benchmarks run eagerly: inference measures a forward pass without gradients,
whereas forward--backward timing includes an output reduction, backpropagation,
and gradient clearing, but no task loss or optimizer update. End-to-end models
use \verb|torch.compile| uniformly. Their training step includes gradient
clearing, a forward pass with cross-entropy loss, backpropagation, and an SGD
update. All inputs and targets are GPU-resident, so data loading and host-to-device
transfers are excluded.

\paragraph{Memory.} \verb|torch.cuda.reset_peak_memory_stats| followed by one
untimed step and \verb|torch.cuda.max_memory_allocated|. This measures
allocator high-water mark, not reserved memory.


\subsection{Environment}

primary measurements used an NVIDIA RTX A6000 (compute capability 8.6, 47.5 GiB), Python 3.12.12, PyTorch 2.6.0+cu124, CUDA 12.4, cuDNN 9.1.0, driver 550.163.01, nvcc 12.4.99, and GCC/G++ 11.4.0. The cross-hardware sweep used an NVIDIA RTX PRO 6000 Blackwell Max-Q (compute capability 12.0, 95.0 GiB), Python 3.12.12, PyTorch 2.9.1+cu130, CUDA 13.0, and cuDNN 9.13.0, driver 595.71.05, nvcc 13.0, and GCC/G++ 11.4.0.
